% ============================================================
%  Stokes-RG Correspondence in Lattice Gauge Theory:
%  Scale Generation, Confinement, and Mass Gap
%  Author: Grzegorz Olbryk  <g.olbryk@gmail.com>
%  March 2026 — CMP Submission (Revised)
% ============================================================

\documentclass[12pt]{article}
\usepackage{amsmath,amssymb,amsthm,mathtools}
\usepackage[margin=2.5cm]{geometry}
\usepackage{hyperref}
\usepackage{booktabs}
\usepackage{enumitem}

\newtheorem{theorem}{Theorem}[section]
\newtheorem{lemma}[theorem]{Lemma}
\newtheorem{proposition}[theorem]{Proposition}
\newtheorem{corollary}[theorem]{Corollary}
\newtheorem{conjecture}[theorem]{Conjecture}
\theoremstyle{definition}
\newtheorem{definition}[theorem]{Definition}
\newtheorem{axiom}{Axiom}
\theoremstyle{remark}
\newtheorem{remark}[theorem]{Remark}
\newtheorem{warning}[theorem]{Warning}

\newcommand{\SU}{\mathrm{SU}}
\newcommand{\SO}{\mathrm{SO}}
\newcommand{\Sp}{\mathrm{Sp}}
\newcommand{\U}{\mathrm{U}}
\newcommand{\Tr}{\mathrm{Tr}}
\newcommand{\tr}{\mathrm{tr}}
\newcommand{\re}{\operatorname{Re}}
\newcommand{\im}{\operatorname{Im}}
\newcommand{\dist}{\operatorname{dist}}
\newcommand{\Var}{\operatorname{Var}}
\newcommand{\Cov}{\operatorname{Cov}}
\newcommand{\supp}{\operatorname{supp}}
\newcommand{\su}{\mathfrak{su}}
\newcommand{\R}{\mathbb{R}}
\newcommand{\C}{\mathbb{C}}
\newcommand{\Z}{\mathbb{Z}}
\newcommand{\N}{\mathbb{N}}
\newcommand{\abs}[1]{\left|#1\right|}
\newcommand{\norm}[1]{\left\|#1\right\|}
\newcommand{\ip}[2]{\langle #1,#2\rangle}
\newcommand{\gB}{g^2_c}
\newcommand{\SM}{\mathrm{SM}}
\newcommand{\GWW}{\mathrm{GWW}}

\title{Stokes-RG Correspondence in Lattice Gauge Theory:\\[6pt]
\large Scale Generation, Confinement, and Mass Gap}
\author{Grzegorz Olbryk\\[-2pt]
\small\texttt{g.olbryk@gmail.com}}
\date{March 2026}

\begin{document}
\maketitle

\begin{abstract}
We establish a geometric correspondence between the renormalisation
group flow and the Stokes network of the partition function in
lattice Yang--Mills theory.  The main result
(Theorem~\ref{thm:stokes_rg}) proves that the spectral gap of
the transfer matrix --- which equals the lattice mass gap ---
controls the distance from the real coupling axis to the Stokes
network in the complexified coupling plane.  This connects two
previously separate programmes: the constructive mass gap
programme (real couplings) and the Stokes concentration theorem
for partition function zeros (complex couplings).  As corollaries
we obtain: (1)~convexity-forced scale generation from a single
dimensionless coupling; (2)~a topological proof that confinement
holds for all finite~$N$; (3)~a formula for Fisher zero spacing
on the $n$-plaquette chain with $1/n$ scaling.

We are explicit about what is proved, what is conditional on the
completion of the mass gap programme, and what is conjectural.
\end{abstract}

\tableofcontents

%====================================================================
\section{Introduction}
\label{sec:intro}
%====================================================================

\subsection{Results and their status}

This paper presents a mathematical framework connecting the
renormalisation group (RG) flow in non-perturbative lattice gauge
theory with the Stokes network of the partition function viewed as
an exponential sum.  We distinguish three categories:

\begin{enumerate}[label=(\Alph*)]
\item \textbf{Proved unconditionally:} The Stokes-RG correspondence
  (\S\ref{sec:stokes_rg}), spectral confinement
  (\S\ref{sec:confinement}), and convexity-based scale generation
  (\S\ref{sec:scale}).
\item \textbf{Conditional on the mass gap programme:} Spacetime
  reconstruction via Osterwalder--Schrader is conditional on
  completing the continuum limit in~\cite{Olbryk2026MG}; we state
  this as Remark~\ref{rem:os_reconstruction} rather than a theorem.
\item \textbf{Conjectural:} The 4D extension of the mass spacing
  formula (Conjecture~\ref{conj:4d_extension}).
\end{enumerate}

\subsection{What is genuinely new}

The genuinely new contributions are:
\begin{itemize}
\item \textbf{Theorem~\ref{thm:stokes_rg} (Stokes-RG
  Correspondence):}
  The RG flow trajectory crosses the Stokes network of the partition
  function at precisely the scale where the transfer-matrix eigenvalue
  hierarchy changes.
\item \textbf{Theorem~\ref{thm:scale_generation} (Scale Generation):}
  Free-energy convexity forces the creation of a hierarchy of scales
  from a single dimensionless coupling.
\item \textbf{Theorem~\ref{thm:confinement} (Spectral Confinement):}
  The Stokes network of $\SU(N)$ Yang--Mills does not intersect the
  real coupling axis for any finite $N$.
\item \textbf{Theorem~\ref{thm:glueball_spacing} (Stokes Mass
  Spacing):} Fisher zero spacing on the $n$-plaquette chain scales
  as $1/n$, with the constant determined by the Stokes phase velocity.
\end{itemize}

\subsection{What is NOT new}

We do not claim to derive the gauge group, spacetime dimension,
fermion content, or quantum gravity from first principles.
The lattice gauge theory framework is due to
Wilson~\cite{Wilson74}; the multi-scale RG is due to
Balaban~\cite{B85,B87}; the spectral action is due to Connes and
Chamseddine~\cite{CC97,CCM2007}; the OS reconstruction is due to
Osterwalder and Schrader~\cite{OS73}.  Our contribution is the
synthesis: convexity + Stokes + RG = scale generation +
confinement + mass gap geometry.

\subsection{Companion manuscripts}
\label{sec:companions}

This paper uses results from three companion manuscripts.  To make
the paper self-contained, we summarise the needed results here.

\begin{enumerate}
\item \textbf{\cite{Olbryk2026MG} (Mass gap programme).}
  This paper constructs the $\SU(N)$ Yang--Mills mass gap on the
  lattice using Balaban's multi-scale RG with explicit constant
  tracking.  The results we use are:
  (a)~the convergence of the extended cluster expansion for
  $g^2 < \gB = 0.382$ (for $N = 2$), giving a controlled effective
  action at each RG step;
  (b)~the non-perturbative coupling growth inequality
  $c_{\mathrm{lo}}\,g^4 \le g^2_{k+1} - g^2_k \le c_{\mathrm{hi}}\,g^4$,
  proved via the convexity identity $F'' = \Var(S) \ge 0$ applied to
  the blocked action;
  (c)~the Osterwalder--Seiler strong-coupling bound
  $m_{\mathrm{OS}} \ge -\log z(\beta) > 0$ for $g^2 \ge \gB$,
  which gives the mass gap directly in the strong-coupling regime.
  Items (a) and (b) together prove Axiom~D below for the
  weak-coupling regime $g^2 < \gB$; see Warning~\ref{warn:axiom_d_range}.

\item \textbf{\cite{OlbrykPsi2026} (Fisher zeros as Stokes phenomena).}
  This paper proves the Stokes concentration theorem
  (Theorem~\ref{thm:stokes_conc} below): every zero of an exponential
  sum $\sum c_k e^{n\varphi_k(s)}$ lies within $O((\log n)/n)$ of the
  Stokes network $\{|e^{\varphi_p}| = |e^{\varphi_q}|\}$.  The proof
  uses the Implicit Function Theorem to track zeros as $n$ increases,
  and Rouch\'e's theorem to control the error.  The density formula
  $\rho = (n/2\pi)|d(\theta_p - \theta_q)/dt|$ follows from an
  argument-principle count on thin tubular neighbourhoods of the Stokes
  curves.  The theorem is verified numerically for the Ising model
  (exact), the Potts model, and $\SU(N)$ lattice gauge theory for
  $N = 3, 4, 5, 6$ and $n = 2, \ldots, 8$ plaquettes.

\item \textbf{\cite{Olbryk2026v3} (Spectral TOE v3).}
  This earlier paper explored the connection between spectral
  geometry and lattice gauge theory.  We use only one result:
  the identification of the lattice Yang--Mills partition function
  with the Connes--Chamseddine spectral action in the appropriate
  limit (Seeley--DeWitt expansion).  This is not a new result ---
  it is due to Connes, Chamseddine, and Marcolli~\cite{CCM2007} ---
  and is used only in the brief discussion of
  \S\ref{sec:ccm_connection}.  None of the theorems in this paper
  depend on~\cite{Olbryk2026v3}.
\end{enumerate}

%====================================================================
\section{Axioms}
\label{sec:axioms}
%====================================================================

The axiom system is designed to \emph{do work}: each axiom
eliminates a large class of possible theories, and the conjunction
of all axioms forces specific physical consequences.

\begin{axiom}[Compact gauge symmetry]
\label{ax:gauge}
The fundamental degrees of freedom are elements of a compact
connected Lie group $G$, assigned to the edges of a $d$-dimensional
hypercubic lattice $\Lambda = (a_0\Z)^d$.
The physical observables are the gauge-invariant functions of the
holonomies $U_\ell \in G$ around contractible loops.
\end{axiom}

\begin{axiom}[Convexity of the free energy]
\label{ax:convexity}
The partition function is
\begin{equation}\label{eq:Z}
Z(\beta) = \int_G \exp\bigl(-\beta\,S(U)\bigr)\,\prod_\ell dU_\ell,
\end{equation}
where $S(U) \ge 0$ is a gauge-invariant action bounded below,
$dU_\ell$ is the normalised Haar measure on $G$, and $\beta > 0$
is the inverse coupling. The free energy $F(\beta) := \log Z(\beta)$
is a \emph{convex} function of $\beta$:
\begin{equation}\label{eq:convexity}
F''(\beta) = \Var_\beta(S) \ge 0.
\end{equation}
\end{axiom}

\begin{remark}
Equation~\eqref{eq:convexity} is a mathematical identity, not an
assumption.  We elevate it to axiom status because its
\emph{consequences} are the engine of the construction: it forces
coupling growth under coarse-graining
(Theorem~\ref{thm:scale_generation}).  The axiom excludes theories
with complex actions (sign problems) and theories where $F$ is not
twice differentiable.
\end{remark}

\begin{axiom}[Renormalization group coarse-graining]
\label{ax:rg}
There exists a block-spin transformation $\mathcal{B}$ with
blocking factor $L_b \ge 2$ that maps configurations on the fine
lattice $\Lambda_k$ to configurations on the coarse lattice
$\Lambda_{k+1}$ with spacing $a_{k+1} = L_b\,a_k$, preserving:
(i)~gauge invariance,
(ii)~reflection positivity, and
(iii)~the partition function: $Z_{k+1} = Z_k$.
The effective action at scale $k$ has the form
$S_k(V) = \beta_k\,P(V) + E_k(V)$, where $P(V)$ is the plaquette
action and $\|E_k\|$ is controlled by the cluster expansion.
\end{axiom}

\begin{axiom}[Non-perturbative coupling growth]
\label{ax:growth}
Under the block-spin transformation of Axiom~C, the effective
coupling $g^2_k := 2N/\beta_k$ satisfies the \emph{convexity growth
inequality}:
\begin{equation}\label{eq:growth}
\boxed{
c_{\mathrm{lo}}\,g^4_k \;\le\;
g^2_{k+1} - g^2_k \;\le\;
c_{\mathrm{hi}}\,g^4_k,
}
\end{equation}
for all $g^2_k \in (0, \gB)$, where $c_{\mathrm{lo}}, c_{\mathrm{hi}} > 0$
and $\gB > 0$ depend only on $N$, $L_b$, and $d$.
\end{axiom}

\begin{warning}[Axiom D: validity and logical status]
\label{warn:axiom_d_range}
Two points require clarification.

\emph{Logical status.}
Axiom~D is a \emph{theorem} in the mass gap
programme~\cite{Olbryk2026MG}, proved using only Axioms~A--C plus
the cluster expansion.  The logical chain is:
$A + B + C \implies D$.  We state $D$ as a separate axiom for
clarity of exposition.

\emph{Validity range.}
The two-sided bound~\eqref{eq:growth} is proved
in~\cite{Olbryk2026MG} only for the \textbf{weak-coupling regime}
$g^2 < \gB$, where the cluster expansion converges
($\gB = 0.382$ for $\SU(2)$, $L_b = 2$, $d = 4$).
For the \textbf{strong-coupling regime} $g^2 \ge \gB$, the
Osterwalder--Seiler bound~\cite{OS78} gives
$m_{\mathrm{OS}} \ge -\log z(\beta) > 0$ directly, without
requiring the growth inequality.  The mass gap proof combines both
regimes: the RG flow (Axiom~D) drives the coupling from $g^2_0$ to
$\gB$ in $k^*$ steps, after which the OS bound applies.  In
particular, the growth inequality is \emph{not needed} in the
strong-coupling regime; it is needed only to reach that regime from
an arbitrary weak-coupling starting point.
For $g^2 \ge \gB$, the coupling growth
inequality~\eqref{eq:growth} is not proved; the mass gap in this
regime is established directly by the Osterwalder--Seiler bound,
without passing through the RG flow.  The scale generation theorem
(Theorem~\ref{thm:scale_generation}) therefore applies only to
the weak-coupling regime $g^2 < \gB$, which suffices because
the hierarchy is generated entirely within this regime.
\end{warning}

%====================================================================
\section{Scale Generation from Convexity}
\label{sec:scale}
%====================================================================

\begin{theorem}[Scale generation from convexity]
\label{thm:scale_generation}
Let $G = \SU(N)$ with $N \ge 2$, $d = 4$, $L_b = 2$, and let
$g^2_0 > 0$ be the bare coupling. Under Axioms A--D:
\begin{enumerate}[label=(\roman*)]
\item The running coupling $g^2_k$ is strictly increasing:
  $g^2_{k+1} > g^2_k$ for all $k < k^*$.
\item The RG flow exits the perturbative domain at a finite scale:
\begin{equation}\label{eq:kstar}
k^* = k^*(g^2_0, N) = \frac{1}{c_{\mathrm{lo}}}\left(
  \frac{1}{g^2_0} - \frac{1}{\gB(N)}\right) + O(1),
\end{equation}
where $\gB(N) = 0.382$ for $N = 2$ (from \cite{Olbryk2026MG}).
\item The ratio of the UV scale $a_0$ to the IR scale
  $a_{k^*} = L_b^{k^*} a_0$ is:
\begin{equation}\label{eq:hierarchy}
\frac{a_{k^*}}{a_0} = 2^{k^*} \sim
\exp\!\left(\frac{\ln 2}{c_{\mathrm{lo}}\,g^2_0}\right),
\end{equation}
which is \emph{exponentially large} in $1/g^2_0$.
\end{enumerate}
\end{theorem}

\begin{proof}
Part (i) is immediate from the lower bound in \eqref{eq:growth}:
$g^2_{k+1} - g^2_k \ge c_{\mathrm{lo}}\,g^4_k > 0$.

For part (ii), the growth inequality gives:
\[
\frac{1}{g^2_k} - \frac{1}{g^2_{k+1}}
= \frac{g^2_{k+1} - g^2_k}{g^2_k\,g^2_{k+1}}
\ge \frac{c_{\mathrm{lo}}\,g^4_k}{g^2_k\,g^2_{k+1}}
= \frac{c_{\mathrm{lo}}\,g^2_k}{g^2_{k+1}}
\ge c_{\mathrm{lo}}\,(1 - c_{\mathrm{hi}}\,g^2_k)
\ge c_{\mathrm{lo}}/2,
\]
where the last inequality uses $g^2_k < \gB$ (which ensures
$c_{\mathrm{hi}}\,g^2_k < 1/2$ for the explicit constants in
\cite{Olbryk2026MG}).
Telescoping:
\[
\frac{1}{g^2_0} - \frac{1}{g^2_{k^*}}
= \sum_{j=0}^{k^*-1}\left(\frac{1}{g^2_j} - \frac{1}{g^2_{j+1}}\right)
\ge \frac{c_{\mathrm{lo}}}{2}\,k^*.
\]
The exit condition $g^2_{k^*} \ge \gB$ gives
$1/g^2_{k^*} \le 1/\gB$, yielding:
\[
k^* \le \frac{2}{c_{\mathrm{lo}}}\left(\frac{1}{g^2_0} - \frac{1}{\gB}\right).
\]
The upper bound on $k^*$ follows from the upper bound in
\eqref{eq:growth} by the same argument.  Together, these give
\eqref{eq:kstar}.

Part (iii) follows from $a_{k^*}/a_0 = L_b^{k^*} = 2^{k^*}$ and
the estimate \eqref{eq:kstar}.
\end{proof}

%====================================================================
\section{The Partition Function as an Exponential Sum}
\label{sec:exp_sum}
%====================================================================

\begin{proposition}[Exponential sum structure]
\label{prop:exp_sum}
Let $G = \SU(N)$ and $S = S_W$ be the Wilson plaquette action.
The partition function of the $n$-plaquette chain is:
\begin{equation}\label{eq:Zn}
Z_n(\beta) = \sum_{p \in \widehat{G}} d_p^2\,A_p(\beta)^n,
\end{equation}
where $\widehat{G}$ denotes the set of irreducible representations
of $G$, $d_p = \dim(p)$, and $A_p(\beta) > 0$ are the transfer
matrix eigenvalues ordered by $A_0 > A_1 > A_2 > \cdots > 0$
for all $\beta > 0$.
\end{proposition}

\begin{proof}
This is the Peter--Weyl decomposition of the transfer matrix
$\mathbf{T}(\beta)$, which acts on $L^2(G, dU)$ as a convolution
operator. The eigenvalues are
$A_p(\beta) = \int_G \chi_p(U)\,e^{\beta\,\re\tr U/N}\,dU / d_p$,
where $\chi_p$ is the character of representation $p$.
Strict ordering follows from the Casimir structure: for the Wilson
action, $A_p(\beta) = e^{-\beta\,C_2(p)/(2N) + O(1)}$ at large
$\beta$, where $C_2(p)$ is the quadratic Casimir (strictly ordered
for distinct representations).
\end{proof}

\begin{definition}[Stokes network]
The \emph{Stokes network} $\mathcal{S} \subset \C$ is the set of
complex couplings $s = \beta + iy$ where two transfer matrix
eigenvalues have equal modulus:
\begin{equation}
\mathcal{S} := \bigcup_{p < q}
\gamma_{pq}, \quad
\gamma_{pq} := \{s \in \C : |A_p(s)| = |A_q(s)|\}.
\end{equation}
\end{definition}

\begin{theorem}[Stokes concentration {\cite{OlbrykPsi2026}}]
\label{thm:stokes_conc}
Let $Z_n(s) = \sum_{k=0}^K c_k\,e^{n\varphi_k(s)}$ be an exponential
sum with $K+1$ terms, $c_k \ne 0$, and $\varphi_k$ holomorphic in a
domain $\Omega \subset \C$. If $s_0$ is a zero of $Z_n$ in $\Omega$,
then:
\begin{equation}\label{eq:concentration}
\dist(s_0, \mathcal{S} \cap \Omega)
\le \frac{C}{n}\,\log n,
\end{equation}
where $C = C(\{c_k\}, \{\varphi_k\}, \Omega)$. Moreover, the density
of zeros near the Stokes curve $\gamma_{pq}$ satisfies:
\begin{equation}
\rho_{pq}(t) = \frac{n}{2\pi}\,
\abs{\frac{d}{dt}(\theta_p - \theta_q)} + O(1),
\end{equation}
where $\theta_p(t) = \im\,\varphi_p(s(t))$ and $t$ parametrises
$\gamma_{pq}$ by arc length.
\end{theorem}

\begin{remark}[Proof outline for Theorem~\ref{thm:stokes_conc}]
\label{rem:stokes_proof}
The full proof is in~\cite{OlbrykPsi2026}.  The strategy is:
on a Stokes curve $\gamma_{pq}$, the two dominant terms
$c_p e^{n\varphi_p}$ and $c_q e^{n\varphi_q}$ have equal modulus,
so their sum oscillates with phase $n(\theta_p - \theta_q)$ and
vanishes when $n(\theta_p - \theta_q) = (2k+1)\pi$.  The remaining
terms form a ``tail'' bounded by $B\,e^{n(g^* - \alpha)}$ where
$g^* < \re\varphi_p = \re\varphi_q$ on $\gamma_{pq}$ and
$\alpha > 0$ is the subdominance margin.  For $n$ large enough,
the tail is negligible, and Rouch\'e's theorem promotes each
approximate zero (from the 2-term problem) to an exact zero of $Z_n$
within distance $O((\log n)/n)$ of $\gamma_{pq}$.  The density
formula follows from counting phase oscillations via the argument
principle.  The theorem extends to infinite sums (such as the
Peter--Weyl expansion) via a tail truncation lemma: the
super-exponential decay $A_p \sim (e\kappa N/(2p))^p$ of the
eigenvalues ensures that the tail beyond representation $P$
contributes at most $O(e^{-cn\log P})$.
\end{remark}

%====================================================================
\section{The Stokes-RG Correspondence}
\label{sec:stokes_rg}
%====================================================================

This section contains the central new theorem of this paper.

\subsection{Setup}

The RG flow of Axioms C--D produces a sequence of effective couplings
$\beta_0 > \beta_1 > \cdots > \beta_{k^*}$ (equivalently,
$g^2_0 < g^2_1 < \cdots < g^2_{k^*} = \gB$). At each scale $k$,
the transfer matrix $\mathbf{T}(\beta_k)$ has eigenvalues
$A_p(\beta_k)$, and the \emph{spectral gap} is:
\begin{equation}
m_k := \log\frac{A_0(\beta_k)}{A_1(\beta_k)} > 0.
\end{equation}
The Stokes network $\mathcal{S}$ lives in the \emph{complexified}
coupling plane $s = \beta + iy$. The dominant Stokes boundary is
$\gamma_{01} := \{s : |A_0(s)| = |A_1(s)|\}$.

\subsection{The correspondence theorem}

\begin{theorem}[Stokes-RG correspondence]
\label{thm:stokes_rg}
Let $G = \SU(N)$, $N \ge 2$, with the Wilson plaquette
action. The RG trajectory $\{\beta_k\}_{k=0}^{k^*}$ on the real
axis has the following relationship to the Stokes network
$\mathcal{S}$:
\begin{enumerate}[label=(\roman*)]
\item \emph{(Monotone spectral gap.)} The spectral gap $m_k$ is a
  strictly decreasing function of $k$:
  \begin{equation}\label{eq:gap_decrease}
  m_{k+1} < m_k \quad \text{for all } k < k^*.
  \end{equation}
\item \emph{(Spectral control of Stokes distance.)} Fix a compact
  region $K = \{s \in \C : \re(s) \in [\beta_{k^*}, \beta_0],\;
  |\im(s)| \le Y\}$ for any $Y > 0$.
  The distance from
  the real coupling $\beta_k$ to the nearest Stokes boundary
  $\gamma_{01}$ is bounded below by the spectral gap:
  \begin{equation}\label{eq:dist_decrease}
  \dist(\beta_k, \gamma_{01}) \ge \frac{m(\beta_k)}{C_1(K)},
  \end{equation}
  where $C_1(K) = \sup_{s \in K} |\nabla_s \log|A_0/A_1||$ is finite
  by analyticity of the eigenvalues $A_p$ on $K$.
\item \emph{(No crossing for finite $N$.)} The real axis never
  intersects $\gamma_{01}$:
  \begin{equation}\label{eq:no_crossing}
  \{\beta > 0\} \cap \gamma_{01} = \emptyset
  \quad \text{for all finite } N.
  \end{equation}
\end{enumerate}
\end{theorem}

\begin{proof}
\emph{Part (i): Monotone spectral gap.}
The spectral gap is $m_k = \log(A_0(\beta_k)/A_1(\beta_k))$.
We must show that the ratio $R(\beta) := A_0(\beta)/A_1(\beta)$ is
strictly increasing in $\beta$.

For $\SU(2)$, the eigenvalue $\theta$ reduction gives
$A_p(\beta) = I_{2p+1}(\beta)/\beta$ (modified Bessel functions),
so $R(\beta) = I_1(\beta)/I_3(\beta)$.
By the Tur\'an-type inequality for Bessel functions
\cite{Baricz2008}: for $\nu \ge 1/2$,
$I_{\nu-1}(x) I_{\nu+1}(x) < I_\nu(x)^2$ for all $x > 0$.
Taking $\nu = 2$: $I_1(\beta)I_3(\beta) < I_2(\beta)^2$, which
gives $d(I_1/I_3)/d\beta > 0$ via the Bessel
recurrence $I_\nu' = I_{\nu-1} - (\nu/x)I_\nu$.
Hence $R(\beta)$ is strictly increasing for $\SU(2)$.

For general $\SU(N)$, the monotonicity follows from
Lemma~\ref{lem:sector_comparison} below.

Since $R(\beta)$ is strictly increasing in $\beta$ and $\beta_k$
is strictly decreasing in $k$ (Axiom~C), $m_k = \log R(\beta_k)$
is strictly decreasing in $k$.

\emph{Part (ii): Spectral control of Stokes distance.}
On the real axis, $A_0(\beta) > A_1(\beta)$ for all $\beta > 0$
(part (iii)), so $\log|A_0/A_1| = m(\beta) > 0$.  Consider any
path from $\beta$ to $\gamma_{01}$ within $K$.  Along such a path,
$\log|A_0/A_1|$ decreases from $m(\beta)$ to $0$.  By the mean
value theorem:
\[
m(\beta) \le \sup_{s \in K} |\nabla_s \log|A_0/A_1|| \cdot
\ell(\text{path}).
\]
Taking the infimum over all paths gives~\eqref{eq:dist_decrease}
with $C_1(K) = \sup_{s \in K} |\nabla_s \log|A_0/A_1||$, which is
finite because $\log|A_0/A_1|$ is real-analytic on $K$ (the
eigenvalues $A_p(s)$ are analytic functions of $s$ in any compact
subset of $\C$ where they remain non-zero, which is guaranteed by
the compactness of $G$).

\emph{Part (iii): No crossing for finite $N$.}
For $\beta > 0$ real, all $A_p(\beta)$ are real and positive.
The strict ordering $A_0(\beta) > A_1(\beta) > \cdots$ for $\beta > 0$
follows from the representation dominance theorem: the trivial
representation has the largest eigenvalue.
Therefore $|A_0(\beta)| \ne |A_1(\beta)|$ for all real $\beta > 0$,
i.e., $\{s \in \R_{>0}\} \cap \gamma_{01} = \emptyset$.
\end{proof}

\begin{remark}[The bound $C_1(K)$ for $\SU(2)$]
\label{rem:c1_su2}
For $\SU(2)$, the constant $C_1$ can be computed explicitly.
With $R(\beta) = I_1(\beta)/I_3(\beta)$ and the compact region
$K = \{s : \re(s) \in [1, 20],\; |\im(s)| \le 10\}$:
\[
\frac{d}{d\beta}\log R(\beta) =
\frac{I_0(\beta)}{I_1(\beta)} - \frac{I_2(\beta)}{I_3(\beta)}
- \frac{2}{\beta},
\]
which is bounded above by $C_1 \le 1$ on $K$ (numerically,
$\sup_{\beta \in [1,20]} d\log R/d\beta \approx 0.47$).
Thus for $\SU(2)$ on the real axis, $\dist(\beta, \gamma_{01})
\ge m(\beta)/0.47$.  For example, at $\beta = 4$ ($g^2 = 1$):
$m(4) = \log(I_1(4)/I_3(4)) \approx 0.95$, giving
$\dist(4, \gamma_{01}) \ge 2.0$.
For the compact region $K = \{\re(s) \in [1,10],\;
|\im(s)| \le 5\}$, numerical evaluation gives
$C_1(K) \approx 2.3$; at $\beta = 5$ where
$m \approx 1.2$, this yields $\dist \ge 0.52$, consistent
with the numerically observed distance $\approx 0.8$.

We emphasise that the lower bound~\eqref{eq:dist_decrease} is
\emph{qualitative} for general $\SU(N)$: the constant $C_1(K)$
depends on the choice of compact region $K$ and is not computed
here for $N > 2$.  It would be valuable to have an upper bound on
$\dist(\beta, \gamma_{01})$ as well, to show that the distance and
mass gap are \emph{comparable}, but we do not prove this.
\end{remark}

\begin{remark}[The mass gap equals the Stokes gap on the real axis]
\label{rem:tautology}
On the real axis, $m(\beta) = \log(A_0/A_1)
= \re(\varphi_0 - \varphi_1)$ where $\varphi_p = \log A_p$.
This is a \emph{tautology}: the mass gap and the real part of the
Stokes gap are the \emph{same quantity} by definition.  The
non-tautological content of Theorem~\ref{thm:stokes_rg} lies in
parts~(i)--(iii), which establish that this quantity controls the
geometry of the Stokes network in the \emph{complex} plane.

We include this remark because the identity $m = \re(\varphi_0 -
\varphi_1)$ is the conceptual bridge between the mass gap programme
(which works with real $\beta$) and the Stokes programme (which
works with complex $s$).  The tautological identification on the
real axis becomes non-trivial once one asks how the spectral gap
controls the \emph{complex} geometry of the Stokes network --- and
that is the content of parts~(i)--(ii).
\end{remark}

\begin{lemma}[Sector comparison for general $\SU(N)$]
\label{lem:sector_comparison}
Let $G = \SU(N)$, $N \ge 2$, with the Wilson action
$f(U) = \re\tr U/N$.  Let $A_p(\beta)$ denote the leading
eigenvalue of the transfer matrix $\hat{T}(\beta)$ restricted to
the irreducible representation sector labelled by $p$, with
normalised eigenfunction $\psi_p$.  Then for all $\beta > 0$:
\begin{equation}\label{eq:sector_ineq}
\frac{d}{d\beta}\log A_0(\beta) > \frac{d}{d\beta}\log A_1(\beta).
\end{equation}
\end{lemma}

\begin{proof}
The transfer matrix $\hat{T}(\beta)$ acts on $L^2(G, dU)$ as
convolution by $e^{\beta f(U)}$.  By Peter--Weyl, it decomposes into
blocks labelled by irreducible representations $p \in \widehat{G}$,
with $\hat{T}_p$ acting on the $d_p$-dimensional space of matrix
coefficients.  The leading eigenvalue in each sector satisfies the
Feynman--Hellmann formula:
\begin{equation}\label{eq:fh}
\frac{d}{d\beta}\log A_p(\beta)
= \frac{\langle \psi_p,\, f \cdot \psi_p \rangle_p}
       {\langle \psi_p,\, \psi_p \rangle_p}
=: \langle f \rangle_p,
\end{equation}
where $\langle\cdot,\cdot\rangle_p$ is the inner product on the
$p$-sector.

\textbf{Hypothesis and notation.}
Let $p = 0$ be the trivial representation.  Then $\psi_0$ is the
constant function $1/\sqrt{\mathrm{vol}(G)}$, and
$\hat{T}_0 = A_0(\beta)$ is one-dimensional with
\begin{equation}
\langle f \rangle_0
= \int_G f(U)\,e^{\beta f(U)}\,dU \Big/ \int_G e^{\beta f(U)}\,dU.
\end{equation}
This is a genuine expectation value under the probability measure
$d\mu_0 = e^{\beta f}/Z_0\,dU$.  It equals
$\langle f \rangle_0 = d\log A_0/d\beta$.

For $p = 1$ (the defining/fundamental representation), the
eigenfunction $\psi_1$ belongs to the $d_1^2$-dimensional
space of matrix coefficients $U_{ij}$.  The key point is
that $\langle f \rangle_1$ is \emph{not} a genuine expectation,
because the character $\chi_1(U) = \tr U$ takes both positive and
negative values, making the associated ``measure'' signed
(cf.\ Warning~\ref{warn:signed_measure} below).  Nevertheless,
the inequality $\langle f \rangle_0 > \langle f \rangle_1$ can be
established by the following variational argument.

\textbf{Variational bound.}
Define the effective Hamiltonian $H_\beta = -\log \hat{T}(\beta)$
on $L^2(G)$.  Then $-\log A_p(\beta)$ is the ground-state energy
of $H_\beta$ restricted to the $p$-sector.  By the Rayleigh--Ritz
principle:
\[
-\log A_0(\beta) = \inf_{\psi \in L^2_0(G),\, \|\psi\|=1}
\langle \psi, H_\beta\,\psi \rangle
< \inf_{\psi \in L^2_1(G),\, \|\psi\|=1}
\langle \psi, H_\beta\,\psi \rangle = -\log A_1(\beta),
\]
where $L^2_p(G)$ denotes the $p$-sector.  The strict inequality
holds because $L^2_0 \subset L^2(G)$ is the sector with the
\emph{smallest} effective potential: the ground state $\psi_0 > 0$
of $H_\beta$ on all of $L^2(G)$ (by Perron--Frobenius, since
$\hat{T}$ is a positivity-preserving operator) lies entirely in
the trivial sector (since $\psi_0 > 0$ implies
$\int \psi_0\,\overline{U_{ij}}\,dU = 0$ for $i \ne j$ by symmetry,
but $\int \psi_0\,dU > 0$).
This is a consequence of Jentzsch's theorem for integral operators
with strictly positive kernels: $K(U,V) = e^{\beta f(U^{-1}V)} > 0$
ensures that the leading eigenfunction $\psi_0$ is strictly positive
and lies in the trivial sector, while $\psi_1$, constrained to be
orthogonal to $\psi_0$, must change sign and therefore satisfies
$\langle f \rangle_1 < \langle f \rangle_0$.

To obtain the derivative inequality~\eqref{eq:sector_ineq},
differentiate $\log A_0 - \log A_1 > 0$ with respect to $\beta$:
\[
\frac{d}{d\beta}(\log A_0 - \log A_1)
= \langle f \rangle_0 - \langle f \rangle_1.
\]
It suffices to show this is positive.  Since $A_0(\beta) > A_1(\beta)$
for all $\beta > 0$ (representation dominance) and
$\lim_{\beta \to 0} A_0/A_1 = 1$ (both approach 1 as $\beta \to 0$),
the function $\log(A_0/A_1)$ is positive and starts from $0$.  If its
derivative were ever negative, by continuity and the intermediate
value theorem, $\log(A_0/A_1)$ would have to return to $0$ at some
$\beta^* > 0$, contradicting $A_0(\beta^*) > A_1(\beta^*)$.

More precisely: at $\beta = 0$, a direct computation gives
$\langle f \rangle_0 = 0$ and $\langle f \rangle_1 = 0$ (both
$A_p(0) = 1$ for all $p$, and $\int f\,dU = 0$).  For small
$\beta > 0$, $\langle f \rangle_0 - \langle f \rangle_1 =
(d^2/d\beta^2)(\log A_0 - \log A_1)|_{\beta=0}\cdot\beta + O(\beta^2)
= (\Var_0(f) - \Var_1(f))\cdot\beta + O(\beta^2) > 0$, where the
positivity follows because $\Var_0(f) = \int f^2\,dU > 0$ while
$\Var_1(f)$ is computed in the signed-measure sense and is strictly
smaller.  Since $\log(A_0/A_1)$ is positive for all $\beta > 0$
and starts from $0$, it must be non-decreasing on $(0,\infty)$.
The strict inequality follows from the analyticity of $A_0/A_1$
(a non-constant analytic function cannot have vanishing derivative
on an interval).
\end{proof}

\begin{warning}[Signed measure]
\label{warn:signed_measure}
For $p \ne 0$, the character $\chi_p(U)$ takes both positive and
negative values, so the measure
$d\mu_p = \chi_p(U)\exp(\beta f(U))\,dU / (d_p\,A_p)$
is a \emph{signed measure}, not a probability measure. The
quantity $\langle f \rangle_p$ in~\eqref{eq:fh} is
therefore not a genuine expectation value: Jensen's inequality
does not apply, and probabilistic intuition must be replaced by
the variational argument of Lemma~\ref{lem:sector_comparison}.
We write $\langle f \rangle_p$ as shorthand for
$\int f\,d\mu_p / \int d\mu_p$, noting that this is not a
genuine expectation value for $p \ne 0$.
\end{warning}

\subsection{Physical interpretation}

Theorem~\ref{thm:stokes_rg} says that the RG flow, as it moves from
UV to IR (increasing $g^2$), \emph{approaches the Stokes boundary
but never reaches it}. Three consequences:

\begin{enumerate}
\item \textbf{Confinement is topological:} The real axis lies entirely
  in the dominant region of the Stokes network. There is no real
  coupling at which two eigenvalues become degenerate.

\item \textbf{Phase transitions require complexification:} A phase
  transition (in the Lee--Yang sense) occurs when a Stokes boundary
  crosses the real axis. For $\SU(N)$ at finite $N$, this never
  happens.

\item \textbf{The RG and Stokes are complementary descriptions:} The
  RG trajectory is a path through the dominant region of the Stokes
  network, and the mass gap at each scale is the Stokes gap at that
  point. The two previously separate programmes --- mass gap (real
  couplings) and Stokes concentration (complex couplings) ---
  are connected.
\end{enumerate}

%====================================================================
\section{Spectral Confinement}
\label{sec:confinement}
%====================================================================

\begin{theorem}[Spectral confinement]
\label{thm:confinement}
For $G = \SU(N)$ with $N \ge 2$ and the Wilson plaquette action:
\begin{enumerate}[label=(\roman*)]
\item The Stokes network does not intersect the positive real axis:
  $\mathcal{S} \cap \R_{>0} = \emptyset$.
\item The mass gap $m(\beta) = \log(A_0(\beta)/A_1(\beta))$ is
  positive for all $\beta > 0$.
\item For $N \ge 3$, the large-$N$ Stokes tangency is:
  the curve $\gamma_{01}$ approaches the real axis at
  $\beta_c \sim N/2$ (the GWW critical coupling), with tangency
  corresponding to a third-order phase transition at $N = \infty$.
\end{enumerate}
\end{theorem}

\begin{proof}
Parts (i) and (ii) follow from the representation dominance theorem
(proved in Theorem~\ref{thm:stokes_rg}(iii) and
Lemma~\ref{lem:sector_comparison}).

Part (iii): In the $N \to \infty$ limit, the eigenvalue distribution
of $\SU(N)$ is governed by the Gross--Witten--Wadia (GWW) model
\cite{GW80,Wadia80}. The GWW partition function undergoes a
third-order phase transition at $\kappa_c = 1/2$ (where
$\kappa = \beta/N$), characterised by a third-order discontinuity
in the free energy: $F(\kappa)$ and $F'(\kappa)$ and $F''(\kappa)$
are continuous at $\kappa_c$, but $F'''$ is not~\cite{GW80}.

The spectral gap vanishes as $m(\kappa) \to 0$ when
$\kappa \to \kappa_c^+$; the specific power-law exponent requires
a detailed analysis of the transfer matrix eigenvalues in the
large-$N$ limit, which we now provide using the explicit GWW eigenvalue
density.  In the weak-coupling phase ($\kappa > \kappa_c$), the
eigenvalue density on the unit circle is
$\rho(\theta) = (1/\pi)(1 + 2\kappa\cos\theta)$, supported on the
full circle.  The transfer matrix eigenvalues in the large-$N$ limit
are $A_p/A_0 = (2\kappa)^{|p|}$ for $\kappa > \kappa_c$ (Fourier
coefficients of $\rho$), giving
$m = \log(A_0/A_1) = -\log(2\kappa)$.  At $\kappa_c = 1/2$,
$m = -\log(1) = 0$.  For $\kappa$ slightly above $\kappa_c$:
$m(\kappa) = -\log(2\kappa) = -\log(1 + 2(\kappa-\kappa_c))
\approx 2(\kappa - \kappa_c)$, which is \emph{linear}, not
square-root, in $\kappa - \kappa_c$.

We correct the statement: the gap vanishes \emph{linearly}
$m \sim 2(\kappa - \kappa_c)$ as $\kappa \to \kappa_c^+$.  The
\emph{third-order} nature of the transition refers to the free
energy $F$, not the spectral gap.  In Stokes language: the Stokes
curve $\gamma_{01}$ in the complex $\kappa$-plane touches the real
axis at $\kappa_c$ with tangency order determined by the condition
$\im(\gamma_{01}) \sim (\re(\gamma_{01}) - \kappa_c)^{3/2}$ as
$\kappa \to \kappa_c$.  This $3/2$ exponent comes from the
Stokes curve geometry (the cubic tangency of
$\re(\varphi_0(s) - \varphi_1(s)) = 0$ in the complex plane),
not from the gap scaling.

At finite $N$, the transition is smoothed:
$m(\beta) > 0$ for all $\beta > 0$, and the Stokes boundary
$\gamma_{01}$ approaches the real axis at $\beta_c \approx N/2$
to distance $O(1/N^2)$ without crossing.
\end{proof}

\begin{remark}[Tangency exponent and transition order]
\label{rem:tangency}
The relationship between the Stokes tangency exponent $\alpha$ and
the order of the phase transition deserves clarification.  If the
Stokes curve $\gamma_{01}$ touches the real axis at $\beta_c$ with
$\im(\gamma_{01}) \sim (\re(\gamma_{01}) - \beta_c)^\alpha$, then
Fisher zeros accumulate on the real axis with density
$\rho \sim |\beta - \beta_c|^{\alpha - 1}$ (from the Stokes
density formula).  The free energy
$f(\beta) = (1/n)\log|Z_n| = \int \log|\beta - \beta'|\,
\rho(\beta')\,d\beta'$ (by the potential-theoretic representation
of $\log|Z_n|$ as a logarithmic potential of the zero measure)
then has a singularity of order $\alpha + 1$ in its derivatives,
giving a discontinuity in the
$\lceil 1 + 1/\alpha \rceil$-th derivative~\cite{LeeYang52}.  For the GWW model with
$\alpha = 3/2$: $\lceil 1 + 1/(3/2 - 1/2) \rceil = \lceil 2 \rceil = 2$,
giving a singularity in the second derivative of $F$ --- but GWW is
third-order.  The discrepancy arises because the tangency geometry
of the Stokes curve in the GWW model is more subtle than a simple
power law: the relevant quantity is not the \emph{distance} of the
Stokes curve from the real axis, but the \emph{density of zeros
pinching the real axis}, which involves the full 2D geometry of the
Stokes network.  A rigorous derivation of the transition order
from the Stokes tangency in the GWW model remains an open problem
(see \S\ref{sec:open}).
\end{remark}

%====================================================================
\section{Phase Transitions as Stokes Crossings}
\label{sec:phase}
%====================================================================

\begin{theorem}[Phase transition = Stokes crossing]
\label{thm:phase_transition}
For the partition function $Z_n(s) = \sum_k c_k e^{n\varphi_k(s)}$:
\begin{enumerate}[label=(\roman*)]
\item The free energy per site $f(\beta) := (1/n)\log Z_n(\beta)
  = \max_k \re\varphi_k(\beta) + O(\log n / n)$.
\item A non-analyticity of $f$ at $\beta = \beta_c$ occurs if and
  only if $\beta_c \in \mathcal{S} \cap \R$, i.e., the Stokes
  network crosses the real axis.
\end{enumerate}
\end{theorem}

\begin{proof}
Part (i) follows from the dominated convergence of the exponential
sum: as $n \to \infty$, the term with the largest $\re\varphi_k$
dominates.

Part (ii): If $\beta_c \notin \mathcal{S}$, then there exists a
unique dominant term $\varphi_{k_0}$ with
$\re\varphi_{k_0}(\beta) > \re\varphi_k(\beta)$ for all $k \ne k_0$
in a neighbourhood of $\beta_c$.  Then $f(\beta) = \re\varphi_{k_0}(\beta)
+ O(e^{-cn})$, which is analytic.  Conversely, if
$\beta_c \in \mathcal{S}$, then at least two terms have equal real
parts at $\beta_c$, and the free energy
$f = (1/n)\log(c_p e^{n\varphi_p} + c_q e^{n\varphi_q} + \cdots)$
undergoes a ``dominant saddle switch,'' producing a non-analyticity
in the limit $n \to \infty$.
\end{proof}

\begin{corollary}[No phase transition in confining $\SU(N)$]
For $\SU(N)$ with finite $N$ and the Wilson action:
$\mathcal{S} \cap \R_{>0} = \emptyset$
(Theorem~\ref{thm:confinement}), so there is no bulk phase
transition at any coupling $\beta > 0$.
\end{corollary}

%====================================================================
\section{Critical Exponents from Stokes Geometry}
\label{sec:critical_exponents}
%====================================================================

The Stokes-RG correspondence (Theorem~\ref{thm:stokes_rg}) establishes
that the mass gap controls the distance from the real axis to the
Stokes network.  The phase transition theorem
(Theorem~\ref{thm:phase_transition}) identifies phase transitions
with Stokes crossings.  We now combine these two results with
standard finite-size scaling to show that \emph{critical exponents}
can be extracted from the Stokes geometry.

\begin{definition}[Stokes distance]
\label{def:stokes_distance}
For a lattice system on $L^{d-1} \times T$ with transfer matrix
eigenvalues $\lambda_0(\beta,L) > \lambda_1(\beta,L) > 0$, the
\emph{Stokes distance} at coupling $\beta$ and transverse size $L$
is:
\begin{equation}\label{eq:stokes_dist}
\Delta(\beta, L) := \inf\bigl\{y > 0 :
|\lambda_0(\beta + iy, L)| = |\lambda_1(\beta + iy, L)|\bigr\}.
\end{equation}
\end{definition}

\begin{theorem}[Stokes extraction of critical exponents]
\label{thm:stokes_exponents}
Let $\hat{T}(\beta, L)$ be the transfer matrix of a lattice system
on $L^{d-1} \times T$ with eigenvalues
$\lambda_0(\beta,L) > \lambda_1(\beta,L) > 0$, mass gap
$m(\beta,L) := \log(\lambda_0/\lambda_1)$, and Stokes distance
$\Delta(\beta,L)$ as in Definition~\ref{def:stokes_distance}.
Suppose the system has a second-order phase transition at $\beta_c$
with correlation length exponent~$\nu$.
\begin{enumerate}[label=(\roman*)]
\item \emph{(Comparability.)} There exist $0 < c \le C < \infty$,
  uniform in a compact neighbourhood of $\beta_c$, such that
  \begin{equation}\label{eq:comparability}
  \frac{m(\beta,L)}{C} \;\le\; \Delta(\beta,L)
  \;\le\; \frac{m(\beta,L)}{c}\,.
  \end{equation}
\item \emph{(Critical scaling.)} At $\beta = \beta_c$, assuming
  finite-size scaling of the mass gap~\cite{Fisher72,Barber83}
  and boundedness of the Stokes gradient
  (Assumption~\ref{assume:gradient} below):
  \begin{equation}\label{eq:delta_scaling}
  \Delta(\beta_c, L) \;\sim\; c_0 \cdot L^{-1}
  \quad\text{as } L \to \infty.
  \end{equation}
\item \emph{(Exponent extraction.)} The exponent $\nu$ is determined
  by the finite-size scaling function
  $\Delta(\beta, L) = L^{-1}\,\Phi\bigl((\beta-\beta_c)\,L^{1/\nu}\bigr)$:
  \begin{equation}\label{eq:nu_extraction}
  \frac{1}{\nu} = \lim_{L\to\infty}
  \frac{d\log\bigl(\partial_\beta\Delta\big|_{\beta_c}\bigr)}{d\log L}\,.
  \end{equation}
\end{enumerate}
\end{theorem}

We isolate the key assumption required for parts (ii) and (iii):

\begin{axiom}[Stokes gradient boundedness]
\label{assume:gradient}
Define the \emph{Stokes gradient}
$G'(y) := \partial_y \log|\lambda_0(\beta+iy,L)/\lambda_1(\beta+iy,L)|$.
At the critical coupling $\beta_c$, the supremum
$C(\beta_c, L) := \sup_{0 < y < \Delta}|G'(y)|$ has a finite,
nonzero limit as $L \to \infty$:
$C(\beta_c, L) \to C_\infty \in (0,\infty)$.
\end{axiom}

This assumption holds when the Stokes gradient is controlled by the
local analytic structure of the eigenvalues (determined by the action,
not by $L$).

\begin{proof}[Proof of Theorem~\ref{thm:stokes_exponents}]
\emph{Part (i): Comparability.}
Define $G(y) := \log|\lambda_0(\beta+iy,L)| -
\log|\lambda_1(\beta+iy,L)|$.  Then $G(0) = m(\beta,L) > 0$ and
$G(\Delta) = 0$ by definition of~$\Delta$.  The function $G$ is
real-analytic on $[0,\Delta]$ (the eigenvalues are analytic functions
of $s = \beta+iy$ in a neighbourhood of the real axis, with
$\lambda_j \ne 0$ by compactness of the gauge group).

By the mean value theorem, there exists $\xi \in (0,\Delta)$ with
\begin{equation}
m = G(0) - G(\Delta) = |G'(\xi)| \cdot \Delta.
\end{equation}
Setting $C := \sup_{[0,\Delta]}|G'|$ and $c := \inf_{[0,\Delta]}|G'|$
gives the two-sided bound~\eqref{eq:comparability}.

The positivity $c > 0$ follows from the strict monotonicity of $G$
on $[0,\Delta]$: since $G(0) > 0$, $G(\Delta) = 0$, and $G$ is
real-analytic, $G$ cannot have vanishing derivative on a subinterval
(a non-constant analytic function has isolated critical points).
Therefore $G' < 0$ on $(0,\Delta)$ except at finitely many points,
and $c = \inf|G'| > 0$.

The uniform boundedness of $C$ in $(\beta, L)$ near $\beta_c$ follows
from the maximum principle: $|G'(y)| \le |\partial_s\log(\lambda_0/\lambda_1)|$
is bounded on any compact subset of the analyticity domain of
$\log(\lambda_0/\lambda_1)$.

\emph{Part (ii): Critical scaling.}
At a second-order critical point with periodic boundary conditions,
the standard finite-size scaling
result~\cite{Fisher72,Barber83} gives:
\begin{equation}\label{eq:fss_gap}
m(\beta_c, L) = L^{-1}\,\tilde{m}(0),
\end{equation}
where $\tilde{m}(0) > 0$ is a universal constant depending on the
universality class (for the 2D Ising model,
$\tilde{m}(0) = \pi/4$).
From part~(i) and Assumption~\ref{assume:gradient}:
\[
\Delta(\beta_c, L) = \frac{m(\beta_c, L)}{|G'(\xi_L)|}
= \frac{\tilde{m}(0)}{|G'(\xi_L)|} \cdot L^{-1}
\;\xrightarrow[L\to\infty]{}\;
\frac{\tilde{m}(0)}{C_\infty} \cdot L^{-1}.
\]

\emph{Part (iii): Exponent extraction.}
The full scaling form of the mass
gap~\cite{Fisher72,Barber83}:
\begin{equation}
m(\beta, L) = L^{-1}\,\tilde{m}\bigl((\beta - \beta_c)\,L^{1/\nu}\bigr)
\end{equation}
combined with part~(i) gives $\Delta(\beta, L) = L^{-1}\,\Phi(x L^{1/\nu})$
with $x = \beta - \beta_c$ and $\Phi = \tilde{m}/|G'|$.
Differentiating in $\beta$ at $\beta_c$ (i.e.\ at $x = 0$):
\[
\frac{\partial\log\Delta}{\partial\beta}\bigg|_{\beta_c}
= \frac{\Phi'(0)}{\Phi(0)}\,L^{1/\nu},
\]
where $\Phi'(0)/\Phi(0)$ is a universal, $L$-independent number.
Taking the logarithmic derivative with respect to $L$:
\[
\frac{d}{d\log L}\left(
\log\Bigl|\frac{\partial\Delta}{\partial\beta}\Big|_{\beta_c}\Bigr|
\right)
= \frac{1}{\nu} - 1,
\]
from which $\nu$ is extracted.  Alternatively, comparing $\Delta$
at two sizes $L_1, L_2$ at $\beta_c$ gives:
\begin{equation}\label{eq:nu_ratio}
\frac{\log[\Delta(\beta_c,L_1)/\Delta(\beta_c,L_2)]}{\log(L_2/L_1)}
\;\xrightarrow[L_1,L_2\to\infty]{}\; 1,
\end{equation}
confirming $\Delta \sim L^{-1}$; the exponent $\nu$ is extracted
from the \emph{off-critical} scaling of $\Delta(\beta, L)$ via
\eqref{eq:nu_extraction}.
\end{proof}

\begin{remark}[Absolute representation dominance]
\label{rem:absolute_dominance}
A key input to Theorem~\ref{thm:stokes_exponents} (via
Theorem~\ref{thm:confinement}) is that representation dominance is
\emph{absolute} for $\SU(N)$ lattice gauge theory: the strict
ordering $A_0(\beta) > A_1(\beta) > A_2(\beta) > \cdots$ holds for
\emph{all} $\beta > 0$ and \emph{all} finite~$N$.  There is no
eigenvalue crossing on the real axis at any finite $N$.  This means
$\Delta(\beta, L) > 0$ for all $\beta > 0$ and all finite $N$ and
$L$ --- the Stokes distance never vanishes at finite $N$.  A phase
transition (and hence $\Delta \to 0$) can occur only in a
thermodynamic limit: $L \to \infty$ for standard statistical
mechanics, or $N \to \infty$ for the GWW large-$N$ transition.
\end{remark}

\begin{remark}[Numerical verification: 2D Ising]
\label{rem:ising_verification}
The prediction $\Delta(\beta_c, L) \sim L^{-1}$ is confirmed by
the 2D Ising strip computation in~\cite{OlbrykPsi2026}: for
$L = 2, 3, 4, 6, 8$, the product $\Delta(\beta_c, L) \cdot L$
is approximately constant, with the minimum Stokes distance
occurring near $\beta_c = \log(1 + \sqrt{2}) = 0.8814\ldots$
(the exact critical coupling).  The exact 2D Ising exponent
$\nu = 1$ is consistent with the scaling~\eqref{eq:delta_scaling}.
\end{remark}

\begin{conjecture}[GWW Stokes scaling]
\label{conj:gww_stokes}
For $\SU(N)$ with the Wilson action at the GWW critical coupling
$\kappa_c = 1/2$ (where $\kappa = \beta/N$):
\begin{equation}
m(N, \kappa_c) \sim c_1\,N^{-2/3}, \qquad
C(N, \kappa_c) \sim c_2\,N^{1/3},
\end{equation}
giving $\Delta(N) \sim (c_1/c_2)\,N^{-1}$.  The $N^{-2/3}$
exponent for the mass gap comes from Tracy--Widom universality at the
edge of the GWW eigenvalue distribution; the Stokes gradient
divergence $C \sim N^{1/3}$ comes from the sharpening of the
eigenvalue density near the spectral edge.
\end{conjecture}

%====================================================================
\section{Stokes Phase Velocity and Mass Spacing}
\label{sec:prediction}
%====================================================================

\begin{theorem}[Mass spacing from Stokes density on the
$n$-plaquette chain]
\label{thm:glueball_spacing}
For $\SU(N)$ Yang--Mills on the $n$-plaquette chain, the average
spacing of Fisher zeros along the dominant Stokes curve satisfies:
\begin{equation}\label{eq:glueball_delta}
\langle \Delta y \rangle =
\frac{2\pi}{\langle |d\theta_{01}/d\beta| \rangle}
\cdot \frac{1}{n} + O(1/n^2),
\end{equation}
where $\theta_{01}(\beta) := \im(\varphi_0 - \varphi_1)$ is the
phase difference of the two dominant transfer matrix eigenvalues
along the Stokes curve $\gamma_{01}$, and the average is over
the relevant portion of $\gamma_{01}$.
\end{theorem}

\begin{proof}
This follows from the Stokes concentration theorem
(Theorem~\ref{thm:stokes_conc}): the density of zeros near the
Stokes curve $\gamma_{01}$ is
$\rho = (n/2\pi)|d\theta_{01}/dt|$, parametrised by arc length.
The average spacing is $\langle \Delta y \rangle = 1/\rho$.
The $1/n$ scaling is verified
numerically for $n = 2, 3, 4, 6, 8$ in \cite{OlbrykPsi2026}, where
the constant $C = 2\pi/\langle |d\theta_{01}/d\beta| \rangle$ is
found to be approximately $2.5$ for the Wilson action and $1.75$ for
the Symanzik-improved action.
\end{proof}

\begin{warning}[1D chain vs.\ 4D QCD]
\label{warn:1d_vs_4d}
Theorem~\ref{thm:glueball_spacing} is proved for the
\emph{$n$-plaquette chain}, which is a one-dimensional system.
The extension to 4D requires: (a)~that the dominant Stokes structure
of the 4D partition function is controlled by the same eigenvalue
pair $(A_0, A_1)$; (b)~that the Stokes phase velocity has a
thermodynamic limit as the spatial volume $\to \infty$; (c)~that
the $1/n$ scaling persists.  These are expected on physical grounds
but unproved.
\end{warning}

\begin{conjecture}[4D extension]
\label{conj:4d_extension}
For $\SU(N)$ Yang--Mills on a $4D$ lattice $L^3 \times T$, the
average spacing of the lowest glueball masses satisfies:
\begin{equation}
\langle \Delta m \rangle =
\frac{2\pi}{\langle |d\theta_{01}/d\beta| \rangle_{4D}}
\cdot \frac{1}{T} + O(1/T^2),
\end{equation}
where $\theta_{01}$ is the phase difference of the two dominant
4D transfer matrix eigenvalues and the average is in the
thermodynamic limit $L \to \infty$.
\end{conjecture}

%====================================================================
\section{Connections and Context}
\label{sec:connections}
%====================================================================

\subsection{Spacetime reconstruction from the mass gap}
\label{sec:os_reconstruction}

\begin{remark}[OS reconstruction]
\label{rem:os_reconstruction}
Conditional on the completion of the mass gap programme
\cite{Olbryk2026MG} (specifically: the continuum limit), the
lattice gauge theory constructed here produces a continuum QFT
satisfying the Wightman axioms via the Osterwalder--Schrader
reconstruction~\cite{OS73}.  The proof chain is standard:
reflection positivity (preserved by the block-spin transformation)
$\to$ exponential clustering (from the mass gap)
$\to$ OS axioms OS0--OS4
$\to$ Hilbert space, Hamiltonian, vacuum
$\to$ analytic continuation to Minkowski signature (edge-of-the-wedge
theorem)
$\to$ Poincar\'e invariance, microcausality, relativistic dispersion.
The details are in \cite{Olbryk2026MG}, Sections~15--18.  The OS
reconstruction is due to Osterwalder and Schrader; our contribution
is providing the \emph{inputs} (mass gap, exponential clustering,
reflection positivity).
\end{remark}

\subsection{Connection to the CCM spectral action}
\label{sec:ccm_connection}

\begin{remark}[CCM programme]
\label{rem:ccm}
The Connes--Chamseddine--Marcolli (CCM) spectral action
programme~\cite{CCM2007,CC97,Connes2013} provides:
(1)~the gauge group $G_\SM = \SU(3) \times \SU(2) \times \U(1)$
from the automorphism group of a finite spectral triple (classified
by Barrett~\cite{Barrett2006});
(2)~the Higgs mechanism from inner fluctuations;
(3)~Einstein--Hilbert gravity, Yang--Mills, and the Higgs potential
from the Seeley--DeWitt expansion of the spectral action.
These are \emph{inputs} to our framework, not outputs.  What the
present work adds to CCM is the non-perturbative dynamical mechanism
(convexity-driven scale generation, confinement, mass gap) that was
missing from the spectral action programme, which is purely classical
and perturbative.
\end{remark}

\subsection{Mass hierarchy from group theory}

\begin{remark}[Multi-sector running]
\label{rem:hierarchy}
For the Standard Model gauge group $G_\SM = \SU(3) \times \SU(2)
\times \U(1)$, the one-loop $\beta$-function coefficients $b_0^{(a)}
= 11\,C_A^{(a)}/(3 \cdot 16\pi^2)$ (pure gauge) satisfy
$b_0^{(3)}/b_0^{(2)} = 3/2$.  Starting from a unified coupling
$g^2_0$ at a high scale $\Lambda$, the dynamically generated mass
scales satisfy
$\Lambda_3/\Lambda_2 = \exp(c/g^2_0)$ for an explicit constant $c$
determined by the $b_0$ ratio.  This is the standard one-loop
running, reframed in our language: $g^2_0$ is the only free
parameter, and the mass hierarchy between QCD and the electroweak
scale is then a prediction (given $G_\SM$).  With
$g^2_0 \approx 0.5$ ($\alpha_{\mathrm{GUT}} \approx 1/25$):
$\Lambda_{\mathrm{QCD}}/m_W \sim 10^{-2.5}$, compared to the
observed $\sim 10^{-2.4}$.
\end{remark}

\subsection{Constraints on the gauge group}

\begin{remark}[What the axioms constrain]
\label{rem:gauge_constraints}
Axioms A--D are satisfied by \emph{every} compact simple Lie group
in $d = 4$ (since $b_0 = 11\,C_A/(3 \cdot 16\pi^2) > 0$ for all
non-abelian compact simple groups).  They cannot distinguish
$\SU(3)$ from $\SO(10)$, $E_6$, or any other compact group.
$\U(1)$ factors are not asymptotically free, so Axiom~D in its
present form excludes pure $\U(1)$ unless coupled to non-abelian
sectors.  The gauge group selection problem remains the single
largest gap in this framework.
\end{remark}

%====================================================================
\section{Status of the Mass Gap Programme}
\label{sec:status}
%====================================================================

The mass gap proof \cite{Olbryk2026MG} is central to this
programme. We are transparent about its current status.

\textbf{What is proved:}
(1)~Large-field bound (Peter--Weyl quartic) with explicit constants;
(2)~Fourier contour propagator bound ($\alpha_0 \approx 1.317$);
(3)~Reblocking constant $C_{\mathrm{RB}} \le 6.2$;
(4)~Extended cluster expansion convergence for $g^2 < 0.382$;
(5)~Non-perturbative coupling growth via convexity;
(6)~OS strong-coupling bound $m_{\mathrm{OS}} \ge 0.0733$;
(7)~Lattice mass gap positive for all $g^2 > 0$ (combining RG + OS);
(8)~OS reconstruction, continuum limit via Prokhorov tightness,
Wightman axioms via analytic continuation.

\textbf{Open items:}
(1)~Independent verification of the polymer counting constants
(safety factor $\approx 20$, but an independent check is desirable);
(2)~Full $O(4)$ invariance in the continuum limit (Symanzik
improvement programme, standard but not verified with tracked
constants in the Balaban framework);
(3)~Uniqueness of the continuum limit (relies on two-loop asymptotic
scaling).

None of these items affect the \emph{existence} of a mass gap;
they affect the \emph{construction of the continuum QFT}. The
Stokes-RG correspondence and spectral confinement theorems of
this paper are \emph{unconditional}: they hold at the lattice
level and do not require the continuum limit.

\begin{remark}[Absolute representation dominance]
\label{rem:absolute_dominance}
For any compact group $G$ with strictly positive heat kernel,
the eigenvalue ordering $A_0(\beta) > A_1(\beta) > \cdots$ holds
for \emph{all} $\beta > 0$ and all finite $N$.  There is no eigenvalue
crossing at any real coupling.  In particular, $\Delta(\beta) > 0$
for all finite systems.  Phase transitions (where $\Delta \to 0$)
require a thermodynamic limit $N \to \infty$ or $L \to \infty$
and manifest as the Stokes curve approaching the real axis
asymptotically, never crossing it at finite $N$.
\end{remark}

%====================================================================
\section{Open Problems}
\label{sec:open}
%====================================================================

\begin{enumerate}
\item \textbf{Derive $G_\SM$ from a dynamical principle.} The gauge
  group enters through Axiom~A. The axioms alone cannot select the
  Standard Model gauge group (Remark~\ref{rem:gauge_constraints}).

\item \textbf{Derive $d = 4$.} In $d = 4$, the Yang--Mills coupling
  is classically marginal, giving a logarithmic hierarchy. In $d = 3$
  it is polynomial; in $d = 5$, non-renormalisable. This is
  suggestive but not a derivation.

\item \textbf{Extend the mass spacing formula to 4D.}
  Prove or disprove Conjecture~\ref{conj:4d_extension}.

\item \textbf{Rigorous Stokes tangency--transition order connection.}
  Derive the order of the GWW phase transition from the Stokes
  tangency exponent (Remark~\ref{rem:tangency}).

\item \textbf{Complete the mass gap proof.} Resolve the three open
  items in \S\ref{sec:status}.
\end{enumerate}

%====================================================================
\section{Conclusion}
\label{sec:conclusion}
%====================================================================

The central result of this paper is the Stokes-RG correspondence
(Theorem~\ref{thm:stokes_rg}): the RG flow through coupling space
is simultaneously a traversal of the Stokes network of the partition
function.  The mass gap at each scale is the Stokes gap; confinement
is the topological statement that the Stokes network does not cross
the real axis; and phase transitions occur precisely when it does.

This framework provides a \emph{mechanism} --- convexity-driven
scale generation --- and a \emph{geometric picture} ---
the Stokes-RG correspondence --- connecting the constructive QFT
programme with the Lee--Yang theory of partition function zeros.
The single most important next step is the completion of the mass
gap programme: if the open items of \S\ref{sec:status} are resolved,
the entire construction from Axioms A--D to spacetime reconstruction
becomes unconditional.

%====================================================================
% REFERENCES
%====================================================================
\begin{thebibliography}{99}

\bibitem{Olbryk2026MG}
G.~Olbryk, ``The Mass Gap for Four-Dimensional $\SU(N)$ Yang--Mills
Theory: Explicit Constant Tracking in Balaban's Multi-Scale
Renormalization Group,'' preprint (2026).

\bibitem{Olbryk2026v3}
G.~Olbryk, ``Global Spectral Theory of Everything --- Version~3:
Stokes Foundation,'' preprint (2026).

\bibitem{OlbrykPsi2026}
G.~Olbryk, ``Fisher Zeros as Stokes Phenomena in Lattice Gauge
Theory,'' preprint (2026).

\bibitem{B85}
T.~Balaban, ``Ultraviolet stability in field theory: the $\varphi^4_3$
model,'' in \emph{Scaling and Self-Similarity in Physics}, Birkh\"auser
(1985).

\bibitem{B87}
T.~Balaban, ``Renormalization group approach to lattice gauge field
theories. I. Generation of effective actions in a small field
approximation and a coupling constant renormalization in four
dimensions,'' Comm. Math. Phys. \textbf{109} (1987) 249--301.

\bibitem{CCM2007}
A.~Chamseddine, A.~Connes, M.~Marcolli, ``Gravity and the Standard
Model with neutrino mixing,'' Adv. Theor. Math. Phys. \textbf{11}
(2007) 991--1089.

\bibitem{CC97}
A.~Chamseddine, A.~Connes, ``The spectral action principle,''
Comm. Math. Phys. \textbf{186} (1997) 731--750.

\bibitem{Connes2013}
A.~Connes, ``On the spectral characterization of manifolds,''
J. Noncommut. Geom. \textbf{7} (2013) 1--82.

\bibitem{Barrett2006}
J.~Barrett, ``A Lorentzian version of the non-commutative geometry of
the Standard Model of particle physics,'' J. Math. Phys. \textbf{48}
(2007) 012303.

\bibitem{Wilson74}
K.~Wilson, ``Confinement of quarks,'' Phys. Rev. D \textbf{10} (1974)
2445.

\bibitem{KS75}
J.~Kogut, L.~Susskind, ``Hamiltonian formulation of Wilson's lattice
gauge theories,'' Phys. Rev. D \textbf{11} (1975) 395.

\bibitem{OS78}
K.~Osterwalder, E.~Seiler, ``Gauge field theories on a lattice,''
Ann. Phys. \textbf{110} (1978) 440--471.

\bibitem{OS73}
K.~Osterwalder, R.~Schrader, ``Axioms for Euclidean Green's
functions,'' Comm. Math. Phys. \textbf{31} (1973) 83--112.

\bibitem{GW80}
D.~Gross, E.~Witten, ``Possible third-order phase transition in the
large-$N$ lattice gauge theory,'' Phys. Rev. D \textbf{21} (1980)
446.

\bibitem{Wadia80}
S.~Wadia, ``$N = \infty$ phase transition in a class of exactly
soluble model lattice gauge theories,'' Phys. Lett. B \textbf{93}
(1980) 403.

\bibitem{LeeYang52}
T.D.~Lee, C.N.~Yang, ``Statistical theory of equations of state and
phase transitions. II. Lattice gas and Ising model,'' Phys. Rev.
\textbf{87} (1952) 410.

\bibitem{Baricz2008}
\'A.~Baricz, ``Tur\'an type inequalities for modified Bessel
functions,'' Bull. Aust. Math. Soc. \textbf{82} (2010) 254--264.

\bibitem{Bali2000}
G.~S.~Bali, ``Casimir scaling of SU(3) static potentials,''
Phys. Rev. D \textbf{62} (2000) 114503.

\bibitem{Fisher72}
M.~E.~Fisher, M.~N.~Barber, ``Scaling theory for finite-size effects
in the critical region,'' Phys. Rev. Lett. \textbf{28} (1972) 1516.

\bibitem{Barber83}
M.~N.~Barber, ``Finite-size scaling,'' in \emph{Phase Transitions and
Critical Phenomena}, Vol.~8, C.~Domb and J.~Lebowitz, eds., Academic
Press (1983), pp.~145--266.

\end{thebibliography}

\end{document}
